% TOP_PROBLEM: {"id": 4800002, "title": "Multiple ergodic averages — Problem 2", "category_id": 11, "category": {"id": 11, "name": "dynamical_systems", "display_name": "Dynamical Systems", "description": "Problems about long-term behavior of deterministic systems, Hamiltonian dynamics, and periodic orbits.", "slug": "dynamical-systems", "order_index": 11, "created_at": "2026-07-31T15:26:25.671Z"}, "status": "partially_solved", "problem_number": "AMR-047-0002", "set_id": 13}
% FIRST_POSTED: 2026-10-01
% ENTRY_KIND: statement_only
% UnsolvedMath ID 4800002; source catalogue code AMR-047-0002
% !TeX program = lualatex
% Definition and sources only; this file is not an LLM solution attempt.
\documentclass[11pt,a4paper]{article}
\usepackage[margin=25mm]{geometry}
\usepackage{fontspec}
\setmainfont{lmroman10-regular.otf}[BoldFont=lmroman10-bold.otf,
 ItalicFont=lmroman10-italic.otf,BoldItalicFont=lmroman10-bolditalic.otf]
\usepackage{amsmath,amssymb,mathtools}
\usepackage{unicode-math}
\setmathfont{latinmodern-math.otf}
\AtBeginDocument{\let\setminus\smallsetminus}
\usepackage{xurl,hyperref}
\hypersetup{colorlinks=true,urlcolor=blue}
\setcounter{secnumdepth}{0}
\setlength{\parindent}{0pt}
\setlength{\parskip}{5pt}
\setlength{\emergencystretch}{3em}
\begin{document}
\section*{Multiple ergodic averages — Problem 2}\label{4800002}
{\small UnsolvedMath ID 4800002 · AMR-047-0002 · Dynamical Systems · AMR Open Problem Lists}
\subsection{Definitions and mathematical statement}
Throughout, a probability-preserving transformation is an invertible
measurable map of a probability space which preserves its measure.
For a bounded measurable function $g$, write $T^ng=g\circ T^n$.
Let $\mathcal C_{T,S}$ be the collection of all complex sequences
$c:\mathbb N\to\mathbb C$ representable as
\[
             c(n)=\int_X f\,(g\circ T^n)\,(h\circ S^n)\,d\mu,
\]
where the probability space $(X,\mathcal B,\mu)$, commuting
probability-preserving transformations $T,S$, and functions
$f,g,h\in L^\infty(\mu)$ may vary.

Define $\mathcal C_T$ by allowing the same representations but
requiring $T=R^a$ and $S=R^b$ for one invertible probability-preserving
transformation $R$ and integers $a,b$. The problem is to prove
$\mathcal C_{T,S}=\mathcal C_T$, or to exhibit a sequence in the
first collection outside the second.

Thus for every particular representation using a commuting pair,
one must find some probability space, one transformation $R$,
integers $a,b$, and three bounded functions whose integrals equal
$c(n)$ for every $n\ge1$. The new system and functions need
not arise as factors or extensions of the original ones.
No ergodicity assumption is imposed in either collection.
Equality means exact equality of sequences, without discarding
a zero-density set of indices or adding an error which vanishes
in Ces\`aro mean. The reverse inclusion follows directly from
commutation of powers; the requested content is the other direction.
\subsection{Short English statement}
Can every double correlation sequence from two commuting transformations be realized exactly using powers of a single transformation?
\subsection{Sources}
\begin{itemize}
\item[S1] ULAM.AI and the UnsolvedMath Contributors, supplied problems.json, record 4800002 (AMR-047-0002), AMR Open Problem Lists. Catalogue identity and source transcription; CC BY 4.0.
\url{https://huggingface.co/datasets/ulamai/UnsolvedMath}.
\item[S2] Frantzikinakis - Some open problems on multiple ergodic averages (2016), Problem 2. Original source indexed by AMR-047-0002.
\url{https://arxiv.org/abs/1103.3808}.
\end{itemize}
\end{document}
